% !TEX root = ../main.tex
% Section 9. Page budget: part of the 4.25 pages for §5-§10 (report/README.md).
% Filled by #75 stage 1 (R7). Sources: PROJECT_CONTEXT.md §10, docs/PROTOCOL.md,
% data/captures/README.md and cap2-cap5 (+ logs, #68), tests/integration/ (#62, #63).
% Frame numbers and counts from tshark 4.2.2 with tools/wireshark/meshsync.lua. Log times are
% defalt's (UTC); kratos logs are 5 h ahead (data/captures/README.md), converted where quoted.
\section{Behaviour under failure (R7)}
\label{sec:failure}

For each failure scenario we give the designed behaviour, what happened on the real hosts, and
the evidence (captures are kratos's view, taken between kratos and defalt; Section~\ref{sec:capture}).

\textbf{Wrong secret} (designed: the HMAC checks reject the peer, ERROR 0x04, close). The secret
is checked twice. First, discovery: kratos's ANNOUNCEs carry a tag made with a different secret,
so defalt drops each one at receive rule~3 and never dials. \code{cap2} has 14 datagrams
(12 ANNOUNCE, 2 LEAVE) and no TCP frame at all; defalt logs
\code{drop D5A6E688 from 192.168.100.92: bad tag} every 5\,s. Second, with a static peer
(\code{--peer}) to reach the TCP handshake (\code{cap2b}): HELLO (frame 6) $\rightarrow$
HELLO\_ACK (8) $\rightarrow$ ERROR \code{\{"code": 4, "msg": "AUTH\_FAILED: bad mac\_r"\}} (10)
$\rightarrow$ FIN from both sides (11, 12). The dialler, defalt, detects the wrong \code{mac\_r}
first (log 13:41:33: \code{initiator handshake failed: AUTH\_FAILED: bad mac\_r}). The attempt
runs four times in the capture, about every 5\,s, and no meshsync message beyond the handshake is
ever sent.
% cap2: 14 frames, 0 tcp; disc types from meshsync.disc.type; cap2-defalt.log 13:38:53-.
% cap2b frames 6/8/10/11/12 and 20/22/24 (11.77 s); 4 x HELLO, 4 x HELLO_ACK, 4 x ERROR;
% cap2b-defalt.log 13:41:33,557. PROTOCOL.md §3.4 rule 3, §5.2, §10 (0x04 fatal).

\textbf{Spoofed or replayed ANNOUNCE} (designed: dropped, bad tag or stale \code{ts}). The bad-tag
case is the one above. The replay check also fired on real traffic: during \code{cap1} defalt
dropped two of kratos's ANNOUNCEs (\code{drop 953B75EA from 192.168.100.92: replay ts 1794662842},
13:27:24). kratos's own timestamps strictly increase in the capture, so defalt most likely
received the same broadcast twice over Wi-Fi, and the check dropped the copy as designed. A
spoofed and a replayed datagram are also tested on localhost
(\code{test\_spoofed\_and\_replayed\_announce\_dropped}).
% cap1-defalt.log (2 "replay ts" lines); data/captures/README.md ("Observations");
% tests/integration/test_failures.py.

\textbf{Untrusted certificate} (designed: TLS alert, close). defalt used a certificate from a
different CA. In \code{cap3}: Client Hello from defalt (frame 7) $\rightarrow$ Server Hello (9)
$\rightarrow$ Alert, level Fatal, description Unknown CA, from defalt (14). The handshake fails
before any meshsync message, and repeats every 5\,s (alerts in frames 14, 27, 40 and 53; no meshsync
message in the whole capture). defalt logs \code{TLS handshake failed (client side):
[SSL: CERTIFICATE\_VERIFY\_FAILED]} (13:55:31) and kratos \code{tlsv1 alert unknown ca}. A TLS/plain
mode mismatch is tested on localhost (\code{test\_tls\_plain\_mode\_mismatch\_closed}).
% cap3: tls.alert_message level 2, desc 48 (Unknown CA) from 192.168.100.80 in frames
% 14/27/40/53; 0 meshsync.type frames; cap3-defalt.log, cap3-kratos.log 18:55:31 (= 13:55:31 UTC).

\textbf{Peer killed mid-transfer} (designed: outstanding chunks are reassigned and the file
completes). We killed defalt, the receiver, with \code{kill -9} during the 100\,MiB transfer. Its
kernel closed the socket at once: a FIN (frame 48679, $t = 91.0$\,s) followed by 421 RSTs for the
data still arriving. kratos logged \code{peer down: defalt ... peer closed connection} (13:58:54
UTC) and discovery expired defalt 11\,s later. kratos's dial attempts at 96.3\,s and 101.4\,s
were refused with RST (frames 49212, 49215), since nothing listened on defalt; after defalt restarted, the handshake ran again (HELLO, frame 49223,
$t = 117.8$\,s) and defalt logged \code{fetching large.bin: 633 of 1403 chunks from 1 peers}: the
770 chunks stored before the kill were kept, not fetched again. The file was then assembled
(\code{assembled large.bin (104857600 bytes, 1403 chunks)}, 13:59:55). With only two hosts in
the capture there was no other source; reassigning chunks to another peer is tested on
localhost with three peers (\code{test\_peer\_killed\_mid\_transfer\_completes\_from\_other\_peer}).
% cap4: tcp.flags.fin frame 48679 (90.988 s) from .80; 421 tcp.flags.reset frames from .80
% before 92 s; SYN from .92 at 96.28/101.41/117.78 s, RST+ACK from .80 in 49212/49215; HELLO (0x01) frame 49223 (117.79 s).
% cap4-kratos.log 18:58:54 / 18:59:05 (expired); cap4-defalt.log 13:59:21,521 and 13:59:55,178.

\textbf{5\,\% packet loss} (designed: completes, slower; TCP retransmissions visible). With
\code{scripts/netem.sh} adding 5\,\% loss on kratos, a 10\,MiB file (144 chunks) still arrived
intact. \code{cap5} holds 12\,417 frames: kratos sent 7\,935 data segments, of which 402 were
retransmissions (5.1\,\%), 52 of them fast retransmissions, and the receiver sent 715 duplicate
ACKs. defalt's fetch took 57.6\,s (14:03:33.3 to 14:04:30.9), about 1.5\,Mbit/s against about
7.9\,Mbit/s without loss in \code{cap1}: TCP kept the data correct and paid in time. Under the same
loss, discovery expired kratos twice while the TCP connection kept working, because broadcast
frames are not retransmitted on Wi-Fi.
% cap5: tcp.analysis.retransmission 402, fast_retransmission 52, duplicate_ack 715;
% ip.src==.92 && tcp.len>0: 7935; total 12417. cap5-defalt.log 14:03:33,325 (fetching),
% 14:04:30,946 (fetched all), 14:04:31,047 (assembled 10485760 bytes), expired 14:04:00, 14:04:55.
% 10485760 x 8 / 57.62 s = 1.46 Mbit/s; cap1 rate from data/captures/README.md ("Observations").

\textbf{Scenarios tested on localhost only.} Concurrent edits on two hosts produce a conflict
copy (\code{test\_concurrent\_edit\_makes\_conflict\_copy}), and a peer that was offline while files
changed catches up through the index exchange when it rejoins
(\code{test\_offline\_peer\_catches\_up}). These run as real \code{meshsync} processes on
\code{127.0.0.1} in \code{tests/integration/} (run in CI), not on the three hosts. The corrupt-chunk
scenario (hash mismatch, re-request from another peer) is designed (\code{PROTOCOL.md} §5.9) but
not yet tested: its test approach is still open.
% tests/integration/test_sync.py, test_failures.py (module docstring: corrupt chunk not there),
% tests/integration/README.md.
